% =========================================================================
% Figura del problema 2.5 (apuntes manuscritos, pág. 27): distribución con
% simetría esférica rho(r) confinada en una esfera de radio R (sombreado
% radial) y geometría del cálculo de los momentos.
% =========================================================================
\begin{figure}[ht]
\centering
\begin{tikzpicture}[x=1cm, y=1cm, >=Stealth]
  \def\Rr{1.7}
  \shade[inner color=magacero!75, outer color=rosaclaro!60] (0,0) circle (\Rr);
  \draw[magacero, line width=0.8pt] (0,0) circle (\Rr);
  \draw[magacero!60, line width=0.4pt, densely dashed] (0,0) ellipse ({\Rr} and {0.42*\Rr});
  \draw[eje] (0,0) -- (0,2.5) node[above, etiq] {$z$};
  \draw[eje] (0,0) -- (2.6,-0.35) node[right, etiq] {$y$};
  \draw[eje] (0,0) -- (-1.3,-1.25) node[below left=-2pt, etiq] {$x$};
  \coordinate (O) at (0,0);
  \coordinate (Zp) at (0,1);
  \coordinate (Q) at (-0.75,1.05);
  \coordinate (P) at (4.6,1.8);
  \draw[vect] (O) -- ($(Q)+(0.06,-0.1)$) node[pos=0.55, left=5pt, etiq, text=white] {$\rp$};
  \cubito{Q}
  \node[etiq, font=\footnotesize, anchor=south east, text=magprof] at ($(Q)+(-0.12,0.12)$) {$\dif q$};
  \draw[vect] (O) -- (P) node[pos=0.62, below right=-2pt, etiq] {$\vc{r}$};
  \draw[vectrojo] ($(Q)+(0.11,0.02)$) -- (P) node[pos=0.62, above=2pt, etiq, text=rojoteo] {$\vc{R}=\vc{r}-\rp$};
  \pic[draw=white, line width=0.5pt, angle radius=0.45cm, "$\theta'$", angle eccentricity=1.75, font=\scriptsize, text=white] {angle=Zp--O--Q};
  \fill[magprof] (P) circle (1.5pt) node[right=2pt, etiq] {$P$};
  \draw[<-, fino] (1.2,-1.2) -- (2.0,-1.75) node[right, etiq, font=\footnotesize] {$\rho=\rho(r')$, \ $r'\le R$};
\end{tikzpicture}
\caption{Problema~2.5. Distribución esféricamente simétrica, $\rho=\rho(r')$, confinada en una esfera de radio $R$ centrada en el origen (el sombreado representa una densidad
decreciente con $r'$). Al usar coordenadas esféricas para el punto fuente, todas las integrales de los momentos se separan en un factor radial y dos angulares, y son estos
últimos los que se anulan.}
\label{fig:sem-esfera-rho}
\end{figure}
